% \begin{document}
\chapter{Introduction}

\dfn{Transition System}{

}
Transition Systems -> basically DFA with possibly infinite states

There can be multiple initial states, and there can be multiple error states. The other states are called "good" $ G \subseteq S \setminus E $.

\dfn{Trace}{
  Alternating sequence of states and labels st at every step the two states and label between is in the relation.
}
A trace always starts from an initial state. A state is reachable if there is a trace that ends in it.

If the set of states if \textbf{finite}, it's possible to compute the reachable states using a graph search algo.

We can abstract a relation without using labels (note that it can become non-deterministic).

\dfn{Composition of relations}{
  We say x and z are \textbf{related}.
}

\dfn{Diagonal relation}{
  Relation to the same state
}

\dfn{Reflexive-transitive closure}{
  \[
  r_1^* = \cup_{n \geq 0}r_1^n
  \]
}

\dfn{Image under relation}{
  \[
    r_1[X] = \{..\}
  \]
}

WE can define a function that 

\dfn{post function}{

}

\thm{Reachable states using post}{
  \[
    \cup_{n \geq 0}\text{post}^n(I) = \text{REach}(M)
  \]
}
\pf{}{
  swap quantifiers
}

In reasoning abt safety

\dfn{Invariant}{
Any superset of reachable states
\[
  \text{Reach}(M) \subseteq P
\]
}
$ P $ is a property satified by all reachable states.

\dfn{Inductive Invariant}{
  An invariant that we can check step-by-step. It contains all initial states and we cannot leave by taking stepps.
}

A way to find that $ P $ is an invariant is to find an inductive invariant that is inside $ P $.

In complicated programms, the set of reachable states can be hard to precicely calculate. So we use Inductive invariants to approximate them.

\mprop{}{
  The smallest inductive invariant is the set of reachable states.
}




% \end{document}
